\section{The Choice at 28}

\subsection{2 Paths and the Ex-Ante Comparison}

Imagine 2 people who are equally capable at 28. Both are engineers. Both are unusually curious, willing to work hard, and capable of learning things that most people find difficult. One accepts a job at a good company. The other starts a company.

10 years later, the first might have a distinguished career, a substantial investment portfolio, a home, and the ability to leave a job without financial panic. The second might have a company valued at 100 million dollars. Or a company that brings in enough cash to employ 12 people and pay its founder handsomely. Or an exhausted bank account and an excellent education in the things that go wrong.

If you want to understand entrepreneurship, you must resist the temptation to compare the employee in the first story with the unusually successful founder in the second. \textbf{The economically relevant comparison is between the choices as they appeared at 28, before anyone knew which outcome would occur.}

\subsection{The Illusion of Retrospective Inevitability}

\looseness=-1 \textbf{In retrospect, successful companies look almost destined to exist.} We explain their triumph with a neat sequence: the founder saw a problem, built something brilliant, found a market, hired exceptional people, and persevered. But at the beginning the founder was surrounded by alternative explanations for everything. Maybe the problem was not important. Maybe customers would not pay. Maybe a competitor would arrive first. Maybe the technology would change. Maybe the founder would lose interest. Maybe another path would have been better.

Entrepreneurship begins with choosing to act amid these competing possibilities. That is what makes it both exciting and hazardous.

\subsection{Wealth, Ambition, and Defining the Game}

People often say they want to start a company because they want to become rich. We should not pretend that this is a shameful or trivial reason. Wealth can buy time, independence, protection from emergencies, the ability to support people one loves, and freedom to pursue difficult work without asking for permission. A sufficiently successful business can also create opportunities for other people on a scale that personal savings cannot.

But \textbf{“I want to be rich” is not yet a theory of entrepreneurship.} It does not tell us whether founding a company is the best way to reach that goal. It does not tell us what kind of company to build, what risks to take, or how to know whether the effort is working. Indeed, it does not even tell us what rich means.

A person earning a large salary and saving consistently may become wealthy. A person who builds a profitable company and owns nearly all of it may become wealthy. A person who creates a fast-growing company may end up holding valuable shares that cannot easily be sold. A person who raises substantial venture capital may control a famous business while owning a comparatively small part of it. Someone else may sell a modest company and obtain more spendable wealth than the founder of a celebrated unicorn.

\textbf{A founder must know what game is he or she planning to play and win.}
